\documentclass{article}
\usepackage{amsmath, amssymb, amsfonts}
\usepackage{geometry}
\usepackage{booktabs}
\usepackage{xcolor}
\geometry{margin=1in}
\title{Orbital Camera Spec vs Implementation Audit}
\author{Rune}
\date{2026-05-11}

\begin{document}
\maketitle

\section{Scope}

Compare \texttt{.opencode/orbital\_camera.tex} (spec) against
\texttt{src/Graphics/Haskan/Camera.hs} (implementation).

\section{Bug 1 (Critical): Pitch rotates around world-X instead of local right axis}

\subsection{Spec (Section 6)}

The spec defines independent angular updates:
\[
\theta_{i+1} = \theta_i + \Delta x \cdot s, \quad
\phi_{i+1} = \text{clip}(\phi_i + \Delta y \cdot s,\; \epsilon,\; \pi - \epsilon)
\]

In spherical coordinates, changing $\theta$ orbits horizontally around $Y$,
and changing $\phi$ moves along the meridian (the great circle through the poles).
The rotation axis for $\phi$ is always the local right vector
$\vec{s} = \text{normalize}(\vec{f} \times \vec{up}_{world})$,
which changes with $\theta$.

\subsection{Implementation}

\begin{verbatim}
yawQ   = axisAngle (V3 0 1 0) yaw      -- world-Y
pitchQ = axisAngle (V3 1 0 0) pitch     -- world-X  <-- WRONG
deltaQ = pitchQ * yawQ
\end{verbatim}

Pitch always rotates around world-$X$, regardless of camera heading.

\subsection{Consequence}

At identity orientation (camera looking along $-Z$), world-$X$ happens to be
the local right axis. Pitch works correctly.

After 90° yaw (camera looking along $-X$), world-$X$ is now the camera's
\textbf{forward} direction. Pitch around world-$X$ produces \textbf{roll}
instead of vertical tilt.

At arbitrary yaw angles, pitch produces a mix of tilt and roll — the apparent
``diagonal inclination'' when the camera is rotated.

\subsection{Fix}

Compute the local right axis from the camera's current heading and yaw
component, then pitch around that:

\begin{verbatim}
(Rotate (V3 yaw pitch _roll)) ->
  let yawQ = axisAngle (V3 0 1 0) yaw
      -- Apply yaw to current orientation to get heading
      headingQ = yawQ * orientation
      forward = normalize (rotate headingQ (V3 0 0 (-1)))
      -- Local right = forward × world-up, in XZ plane
      right = normalize (forward `cross` V3 0 1 0)
      pitchQ = axisAngle right pitch
      newOrientation = pitchQ * headingQ
  in cam { orientation = newOrientation }
\end{verbatim}

\section{Bug 2: Elevation bounds never enforced}

\subsection{Spec (Section 2)}

\[
\phi \in [\epsilon,\; \pi - \epsilon]
\]

\subsection{Implementation}

The \texttt{elevationBounds} field exists but is \textbf{never checked} in
\texttt{orbitalModify}. The \texttt{Rotate} handler applies unrestricted pitch:

\begin{verbatim}
newOrientation = deltaQ * orientation
in cam {orientation = newOrientation}
\end{verbatim}

\subsection{Consequence}

The camera can pass through the poles ($\phi = 0$ or $\phi = \pi$), where the
forward vector becomes parallel to world-up. At that point, linear's
\texttt{lookAt} produces a degenerate view matrix (the cross product
$\vec{f} \times \vec{up}$ collapses to zero), causing the view to snap or
flip.

\subsection{Fix}

After computing \texttt{newOrientation}, extract the elevation angle and clamp:

\begin{verbatim}
let newEl = quatToElevation newOrientation
    (V2 elMin elMax) = fromMaybe (V2 0.01 (pi - 0.01)) elevationBounds
    clampedEl = clamp newEl elMin elMax
    -- Reconstruct orientation with clamped elevation
    az = quatToAzimuth newOrientation
    azQ = axisAngle (V3 0 1 0) az
    elQ = axisAngle (V3 1 0 0) clampedEl
    finalOrientation = elQ * azQ
in cam { orientation = finalOrientation }
\end{verbatim}

Note: this re-decomposes into Euler angles for clamping, which re-introduces
potential gimbal lock at the boundaries. The clamp bounds should be generous
($\epsilon \approx 0.01$ rad) to keep the quaternion approach valid in the
operating range.

\section{Bug 3: Spec's projection matrix is internally inconsistent}

\subsection{The problem}

The spec defines two matrices:

\textbf{View matrix (Section 4):} third row is $-\vec{f}$, producing
right-handed view space where visible objects have $z_\text{view} < 0$.

\textbf{Projection (Section 5):} has $[3][2] = 1$, meaning
$w_\text{clip} = z_\text{view}$.

For visible objects: $z_\text{view} < 0 \Rightarrow w_\text{clip} < 0$.
Vulkan clips vertices with $w < 0$. \textbf{All geometry is clipped.}

\subsection{Implementation}

The implementation correctly uses linear's \texttt{perspective} which has
$[3][2] = -1$, giving $w_\text{clip} = -z_\text{view} > 0$ for visible
objects. The Y-flip is applied separately. This is correct.

\subsection{Fix (spec only)}

Change the spec's projection to:
\[
P = \begin{bmatrix}
\frac{1}{a \cdot \tan(\frac{fov}{2})} & 0 & 0 & 0 \\
0 & -\frac{1}{\tan(\frac{fov}{2})} & 0 & 0 \\
0 & 0 & \frac{-(n+f)}{f-n} & \frac{2nf}{f-n} \\
0 & 0 & -1 & 0
\end{bmatrix}
\]

Or use the [0,1] depth formulation but with $[3][2] = -1$.

\section{Non-Issue: Position parameterization}

The spec uses spherical coordinates:
\begin{align*}
P_x &= r \sin\phi \cos\theta + T_x \\
P_y &= r \cos\phi + T_y \\
P_z &= r \sin\phi \sin\theta + T_z
\end{align*}

The implementation uses quaternion rotation of $(0, 0, r)$:
\begin{verbatim}
target + rotate orientation (V3 0 0 distance)
\end{verbatim}

These produce the same set of positions (the sphere of radius $r$ around $T$)
with different parameterizations. The implementation's identity orientation
$(0, 0, r)$ corresponds to the spec's $\theta = \pi/2, \phi = \pi/2$.
Both are valid; the quaternion approach avoids gimbal lock.

\section{Non-Issue: Forward direction}

\begin{verbatim}
orbitalCameraForward = normalize $ rotate orientation (V3 0 0 (-1))
\end{verbatim}

This equals $\text{normalize}(T - P)$ since:
\[
T - P = -\text{rotate}(q)(0, 0, r) = \text{rotate}(q)(0, 0, -r)
\]
\[
\text{normalize}(T - P) = \text{rotate}(q)(0, 0, -1) \quad \checkmark
\]

\section{Non-Issue: View matrix}

Uses linear's \texttt{lookAt} which produces the exact matrix from spec
Section 4. $\checkmark$

\section{Summary}

\begin{table}[h]
\centering
\begin{tabular}{llp{5.5cm}}
\toprule
\textbf{\#} & \textbf{Severity} & \textbf{Issue} \\
\midrule
1 & Critical & Pitch around world-X instead of local right axis \\
2 & Medium & Elevation bounds (\texttt{elevationBounds}) never enforced \\
3 & Spec-only & Projection $[3][2]=1$ incompatible with right-handed view \\
\bottomrule
\end{tabular}
\end{table}

Bug 1 is the most likely cause of the ``diagonal inclination'' observed when
the camera is rotated away from the default heading.

\end{document}
